% Options for packages loaded elsewhere
\PassOptionsToPackage{unicode}{hyperref}
\PassOptionsToPackage{hyphens}{url}
\documentclass[
  ignorenonframetext,
  aspectratio=169,
]{beamer}
\newif\ifbibliography
\usepackage{pgfpages}
\setbeamertemplate{caption}[numbered]
\setbeamertemplate{caption label separator}{: }
\setbeamercolor{caption name}{fg=normal text.fg}
\beamertemplatenavigationsymbolsempty
% remove section numbering
\setbeamertemplate{part page}{
  \centering
  \begin{beamercolorbox}[sep=16pt,center]{part title}
    \usebeamerfont{part title}\insertpart\par
  \end{beamercolorbox}
}
\setbeamertemplate{section page}{
  \centering
  \begin{beamercolorbox}[sep=12pt,center]{section title}
    \usebeamerfont{section title}\insertsection\par
  \end{beamercolorbox}
}
\setbeamertemplate{subsection page}{
  \centering
  \begin{beamercolorbox}[sep=8pt,center]{subsection title}
    \usebeamerfont{subsection title}\insertsubsection\par
  \end{beamercolorbox}
}
% Prevent slide breaks in the middle of a paragraph
\widowpenalties 1 10000
\raggedbottom
\AtBeginPart{
  \frame{\partpage}
}
\AtBeginSection{
  \ifbibliography
  \else
    \frame{\sectionpage}
  \fi
}
\AtBeginSubsection{
  \frame{\subsectionpage}
}
\usepackage{iftex}
\ifPDFTeX
  \usepackage[T1]{fontenc}
  \usepackage[utf8]{inputenc}
  \usepackage{textcomp} % provide euro and other symbols
\else % if luatex or xetex
  \usepackage{unicode-math} % this also loads fontspec
  \defaultfontfeatures{Scale=MatchLowercase}
  \defaultfontfeatures[\rmfamily]{Ligatures=TeX,Scale=1}
\fi
\usepackage{lmodern}
\ifPDFTeX\else
  % xetex/luatex font selection
\fi
% Use upquote if available, for straight quotes in verbatim environments
\IfFileExists{upquote.sty}{\usepackage{upquote}}{}
\IfFileExists{microtype.sty}{% use microtype if available
  \usepackage[]{microtype}
  \UseMicrotypeSet[protrusion]{basicmath} % disable protrusion for tt fonts
}{}
\makeatletter
\@ifundefined{KOMAClassName}{% if non-KOMA class
  \IfFileExists{parskip.sty}{%
    \usepackage{parskip}
  }{% else
    \setlength{\parindent}{0pt}
    \setlength{\parskip}{6pt plus 2pt minus 1pt}}
}{% if KOMA class
  \KOMAoptions{parskip=half}}
\makeatother
\usepackage{longtable,booktabs,array}
\usepackage{caption}
\captionsetup[table]{skip=6pt}
\newcounter{none} % for unnumbered tables
\usepackage{calc} % for calculating minipage widths
\usepackage{caption}
% Make caption package work with longtable
\makeatletter
\def\fnum@table{\tablename~\thetable}
\makeatother
\setlength{\emergencystretch}{3em} % prevent overfull lines
\providecommand{\tightlist}{%
  \setlength{\itemsep}{0pt}\setlength{\parskip}{0pt}}
% Auto-generated by infrastructure/rendering/_pdf_latex_helpers.py::extract_math_font_preamble.
% Minimal header passed to Pandoc via -H so Beamer slides pick up the same math font
% as the combined-PDF path. Adjust the manuscript's preamble.md to change the font globally.
\usepackage{unicode-math}
\setmathfont{latinmodern-math.otf}

% Manuscript macro/package fallbacks (newcommand -> providecommand).
% Core mathematics
\usepackage{amsmath}
\usepackage{amssymb}
% Document layout
% Tables
\usepackage{booktabs}
% Caption styling (booktabs already loaded above). Small caption font with a
% bold label ("Figure 1", "Table 2") gives the math/figure-heavy layout a
% consistent, journal-like caption rhythm without fighting pandoc-crossref —
% pandoc emits standard \caption{}, and caption/captionsetup only restyle it.
% ── Theorem environments ─────────────────────────────────────────────────
% amsthm is loaded above. The FedGVI<->Friston bridge is stated as a small
% number of theorems/definitions/lemmas/propositions/corollaries; sharing one
% counter (definition/lemma/proposition/corollary all step the `theorem`
% counter) gives a single monotone numbering 1, 2, 3, ... across all kinds, so
% authors can reference them as Theorem 1, Definition 2, Corollary 3 without a
% per-kind counter drifting out of order. Plain style for theorem-like results,
% definition style (upright body) for definitions.
% (algorithm2e intentionally not loaded: this manuscript states results as
% numbered amsthm Theorems/Definitions, not algorithm2e pseudocode.)
% Code listings
\usepackage{listings}
% Typography and formatting
% (siunitx intentionally not loaded: no \SI/\num units appear in this manuscript.)
% Cross-references and citations
% (cleveref and titlesec intentionally not loaded: unavailable in this TeX tree
% and not required. Cross-references resolve through pandoc-crossref's own
% \ref-style names — no \cref appears — and section heading styling falls back
% to the pandoc/LaTeX defaults.)
% ── TeX Live Basic-compatible mono font for code listings ────────────
% Latin Modern Mono is bundled with TeX Live Basic and keeps the manuscript
% renderable on minimal TeX installations. Mathematical Unicode coverage is
% provided separately by Latin Modern Math below, so the code font does not
% need a private JuliaMono installation.
%
% Use the TeX font filename rather than the system-family name: the minimal
% TeX Live fontconfig database may not register the OpenType family even though
% kpsewhich can resolve the bundled files.
% Math font for unicode-math: Latin Modern Math (TeX Live) has full BMP
% coverage including U+2223 (\mid), U+226A/226B (\ll/\gg), and the Greek/
% blackboard letters used in equations. Without an explicit \setmathfont,
% unicode-math falls back to lmroman text font which lacks several glyphs
% and emits "Missing character" warnings on every \mid in math mode.

% Slide layout defaults for warning-clean scientific decks.
\usepackage{etoolbox}
\IfFileExists{xurl.sty}{\usepackage{xurl}}{}
\IfFileExists{seqsplit.sty}{\usepackage{seqsplit}}{\newcommand{\seqsplit}[1]{#1}}
\protected\def\breakseq#1{\seqsplit{#1}}
\protected\def\breaktt#1{\begingroup\ttfamily\seqsplit{#1}\endgroup}
\setlength{\emergencystretch}{6em}
\tolerance=5000
\hbadness=10000
\hfuzz=1pt
\setlength{\tabcolsep}{2pt}
\AtBeginEnvironment{longtable}{\tiny\renewcommand{\arraystretch}{0.86}\setlength{\tabcolsep}{1pt}}
\AtBeginEnvironment{tabular}{\tiny\renewcommand{\arraystretch}{0.86}\setlength{\tabcolsep}{1pt}}
\AtBeginEnvironment{equation}{\tiny}
\AtBeginEnvironment{equation*}{\tiny}
\AtBeginEnvironment{align}{\tiny}
\AtBeginEnvironment{align*}{\tiny}
\AtBeginEnvironment{itemize}{\footnotesize}
\AtBeginEnvironment{enumerate}{\footnotesize}
\AtBeginEnvironment{description}{\footnotesize}
\setbeamerfont{caption}{size=\tiny}
\setbeamerfont{caption name}{size=\tiny}
\setbeamerfont{normal text}{size=\small}
\setbeamerfont{frametitle}{size=\small}
\setbeamerfont{section title}{size=\footnotesize}
\setbeamerfont{subsection title}{size=\footnotesize}
\setbeamertemplate{section page}{%
  \centering
  \begin{beamercolorbox}[sep=12pt,center,wd=\paperwidth]{section title}
    \parbox{0.86\paperwidth}{\centering\usebeamerfont{section title}\insertsection\par}
  \end{beamercolorbox}
}
\setbeamertemplate{subsection page}{%
  \centering
  \begin{beamercolorbox}[sep=8pt,center,wd=\paperwidth]{subsection title}
    \parbox{0.86\paperwidth}{\centering\usebeamerfont{subsection title}\insertsubsection\par}
  \end{beamercolorbox}
}
\setlength{\abovecaptionskip}{2pt}
\setlength{\belowcaptionskip}{0pt}

% Natbib and cross-reference fallbacks — slides don't load natbib
% or cleveref, but combined-PDF manuscript prose may emit these
% commands. The fallback renders citations as a bracketed key list
% and cross-references as detokenized labels so slides don't fail on
% undefined control sequences or raw underscores. \providecommand is
% a no-op if packages load later.
\providecommand{\citep}[1]{[#1]}
\providecommand{\citet}[1]{#1}
\providecommand{\citealp}[1]{#1}
\providecommand{\citeauthor}[1]{#1}
\providecommand{\citeyear}[1]{#1}
\providecommand{\cref}[1]{\texttt{\detokenize{#1}}}
\providecommand{\Cref}[1]{\texttt{\detokenize{#1}}}
\providecommand{\crefrange}[2]{\texttt{\detokenize{#1}}--\texttt{\detokenize{#2}}}
\providecommand{\Crefrange}[2]{\texttt{\detokenize{#1}}--\texttt{\detokenize{#2}}}
\providecommand{\paragraph}[1]{\textbf{#1}\ }

% Opt-in accessible presentation profile.
\setbeamerfont{normal text}{size*={20pt}{24pt}}
\setbeamerfont{frametitle}{size*={28pt}{32pt}}
\setbeamerfont{section title}{size*={28pt}{32pt}}
\setbeamerfont{subsection title}{size*={28pt}{32pt}}
\setbeamerfont{caption}{size*={16pt}{19pt}}
\setbeamerfont{caption name}{size*={16pt}{19pt}}
\setbeamerfont{itemize/enumerate body}{size*={20pt}{24pt}}
\setbeamerfont{itemize/enumerate subbody}{size*={20pt}{24pt}}
\setbeamerfont{itemize/enumerate subsubbody}{size*={20pt}{24pt}}
\setbeamerfont{itemize item}{size*={20pt}{24pt}}
\setbeamerfont{itemize subitem}{size*={20pt}{24pt}}
\setbeamerfont{itemize subsubitem}{size*={20pt}{24pt}}
\setbeamerfont{description body}{size*={20pt}{24pt}}
\setbeamerfont{description item}{size*={20pt}{24pt}}
\setbeamerfont{quote}{size*={20pt}{24pt}}
\setbeamerfont{quotation}{size*={20pt}{24pt}}
\AtBeginEnvironment{longtable}{\fontsize{20pt}{24pt}\selectfont}
\AtBeginEnvironment{tabular}{\fontsize{20pt}{24pt}\selectfont}
\AtBeginEnvironment{equation}{\fontsize{20pt}{24pt}\selectfont}
\AtBeginEnvironment{equation*}{\fontsize{20pt}{24pt}\selectfont}
\AtBeginEnvironment{align}{\fontsize{20pt}{24pt}\selectfont}
\AtBeginEnvironment{align*}{\fontsize{20pt}{24pt}\selectfont}
\AtBeginEnvironment{itemize}{\fontsize{20pt}{24pt}\selectfont}
\AtBeginEnvironment{enumerate}{\fontsize{20pt}{24pt}\selectfont}
\AtBeginEnvironment{description}{\fontsize{20pt}{24pt}\selectfont}
\AtBeginDocument{\usebeamerfont{normal text}}
\newlength{\AFSlideTableWidth}
% The fixed footer reservation is calibrated with the semantic seven-line regular-body budget.
\setbeamertemplate{footline}{%
  \leavevmode\vbox{%
    \hbox{%
      \begin{beamercolorbox}[wd=\paperwidth,ht=16pt,dp=3pt,center]{author in head/foot}%
        \fontsize{16pt}{19pt}\selectfont Untagged PDF derivative \textbar\ \href{../web/index.html}{HTML reader}%
      \end{beamercolorbox}%
    }%
    \vskip6pt%
  }%
}
% Accessible footer reserved height: 25pt.

\newtheorem{proposition}{Proposition}
\newtheorem{hypothesis}{Hypothesis}
\newtheorem{remark}[theorem]{Remark}
\newtheorem{axiom}{Axiom}
\newtheorem{property}{Property}
\usepackage{bookmark}
\IfFileExists{xurl.sty}{\usepackage{xurl}}{} % add URL line breaks if available
\urlstyle{same}
% fallback for those not using the hyperref driver hyperxmp:
\makeatletter
\@ifundefined{xmpquote}{\newcommand{\xmpquote}[1]{#1}}{}
\makeatother
\hypersetup{
  hidelinks,
  pdfcreator={LaTeX via pandoc}}

\author{\texorpdfstring{}{}}
\date{}

\begin{document}

\begin{frame}{Supplement: N-level hierarchical POMDP methods}
\protect\phantomsection\label{sec:supp-3level}
This supplement specifies the generic \(N\)-level architecture that
sec.~14.11 exercises at depth three: how the
meta-context (L3), context (L2), and location (L1) factors are chained
through conditioned priors,
\end{frame}

\begin{frame}[fragile]{Supplement: N-level\ldots{} (part 2)}
what the declarative \texttt{LayerSpec} interface fixes versus leaves
free, and the top-down/bottom-up passes the inference runs.
\end{frame}

\begin{frame}{Supplement: N-level\ldots{} (part 3)}
It answers \emph{what the executed 3-level result is a special case of}
--- the reason the same code runs at other depths without new
mathematics --- while recording that only the declared 3-level
configuration is empirically evaluated here.
\end{frame}

\begin{frame}[fragile]{Generative model for an N-level hierarchy}
\protect\phantomsection\label{generative-model-for-an-n-level-hierarchy}
The 3-level POMDP is built by \texttt{build\_3level\_world} in the
\texttt{fedference.pomdp} module. It extends the 2-level construction
(sec.~14.10) by adding a top-level meta-context
factor.
\end{frame}

\begin{frame}[fragile]{Generative model for an\ldots{} (part 2)}
\textbf{L3 (meta-context).} 2 \texttt{low\_threat} /
\texttt{high\_threat} states begin from a uniform prior and gate the L2
context prior.

\textbf{L2 (context).} 2 \texttt{quiet} / \texttt{alert} states gate the
L1 prior through context-conditioned location priors.
\end{frame}

\begin{frame}{Generative model for an\ldots{} (part 3)}
\textbf{L1 (location).} The standard 3x3 grid has 9 states and sensor
acuity 0.85.

The conditioned priors are (see eq.~38 and
eq.~39):
\end{frame}

\begin{frame}[fragile]{Generative model for an\ldots{} (part 4)}
{\def\LTcaptype{none} % do not increment counter
% template-accessible-table-inset
\begingroup
\setlength{\AFSlideTableWidth}{\textwidth}%
\setlength{\LTleft}{\fill}%
\setlength{\LTright}{\fill}%
\begin{longtable}[]{@{}
  >{\raggedright\arraybackslash}p{(\AFSlideTableWidth - 2\tabcolsep) * \real{0.5172}}
  >{\raggedright\arraybackslash}p{(\AFSlideTableWidth - 2\tabcolsep) * \real{0.4828}}@{}}
\toprule\noalign{}
\begin{minipage}[b]{\linewidth}\raggedright
L3 state
\end{minipage} & \begin{minipage}[b]{\linewidth}\raggedright
L2 prior (quiet, alert)
\end{minipage} \\
\midrule\noalign{}
\endhead
\texttt{low\_threat} & (0.50, 0.50) --- uniform context \\
\texttt{high\_threat} & (0.20, 0.80) --- peaked at alert \\
\bottomrule\noalign{}
\end{longtable}
\endgroup
}
\end{frame}

\begin{frame}[fragile]{Generative model for an\ldots{} (part 5)}
{\def\LTcaptype{none} % do not increment counter
% template-accessible-table-inset
\begingroup
\setlength{\AFSlideTableWidth}{\textwidth}%
\setlength{\LTleft}{\fill}%
\setlength{\LTright}{\fill}%
\begin{longtable}[]{@{}
  >{\raggedright\arraybackslash}p{(\AFSlideTableWidth - 2\tabcolsep) * \real{0.2963}}
  >{\raggedright\arraybackslash}p{(\AFSlideTableWidth - 2\tabcolsep) * \real{0.7037}}@{}}
\toprule\noalign{}
\begin{minipage}[b]{\linewidth}\raggedright
L2 state
\end{minipage} & \begin{minipage}[b]{\linewidth}\raggedright
L1 prior
\end{minipage} \\
\midrule\noalign{}
\endhead
\texttt{quiet} & uniform over all 9 location states \\
\texttt{alert} & mass 0.60 at center cell (flat index 4), residual
uniform \\
\bottomrule\noalign{}
\end{longtable}
\endgroup
}
\end{frame}

\begin{frame}[fragile]{Generic N-level architecture}
\protect\phantomsection\label{generic-n-level-architecture}
The \texttt{LayerSpec} type and \texttt{build\_nlevel\_world} function
in the \texttt{fedference.pomdp} module implement the generic N-level
version. The declarative layer specification is
\texttt{hierarchical\_layers.yaml} under
\texttt{src/fedference/config/};
\end{frame}

\begin{frame}[fragile]{Generic N-level architecture (part 2)}
it mirrors the canonical 3-level defaults (a standalone documentation
artifact not read by any code path, kept in sync with the
\texttt{build\_3level\_world} defaults).
\end{frame}

\begin{frame}[fragile]{Generic N-level architecture (part 3)}
The constructor accepts depth \texorpdfstring{\ensuremath{\ge}}{>=} 2; the executed empirical result in this
manuscript is restricted to the declared 3-level configuration, and the
leaf layer must carry \texttt{n\_states\ ==\ N\_LOCATIONS}.
\end{frame}

\begin{frame}[fragile]{Inference algorithm across hierarchy levels}
\protect\phantomsection\label{inference-algorithm-across-hierarchy-levels}
\texttt{fedference.pomdp.nlevel\_infer} performs 4 passes of top-down /
bottom-up alternating minimization over all N levels:
\end{frame}

\begin{frame}{Inference algorithm across\ldots{} (part 2)}
\textbf{1. Top-down pass.} Compute the empirical prior for each level by
marginalizing over the level above (eq.~38,
eq.~39).
\end{frame}

\begin{frame}{Inference algorithm across\ldots{} (part 3)}
\textbf{2. L1 update.} The one-step variational posterior on the
observation is
\(q_{\text{loc}} = \operatorname{softmax}(\log \widetilde{\pi}_{0,\mathrm{L1}} +
\log A[\text{obs},\,\cdot])\).
\end{frame}

\begin{frame}{Inference algorithm across\ldots{} (part 4)}
\textbf{3. Bottom-up pass.} Update each non-leaf level's belief from the
marginal evidence contributed by the level below:
\(\ell_j = \log(\tilde{p}_{\text{child|parent=}j}^\top q_{\text{child}})\).
\end{frame}

\begin{frame}{Inference algorithm across\ldots{} (part 5)}
After 4 iterations the agent broadcasts all N level beliefs; the colony
federates each level independently via a log-linear pool
(eq.~9).
\end{frame}

\begin{frame}{Study parameters for the three-level run}
\protect\phantomsection\label{study-parameters-for-the-three-level-run}
% template-accessible-table-inset
\begingroup
\setlength{\AFSlideTableWidth}{\textwidth}%
\setlength{\LTleft}{\fill}%
\setlength{\LTright}{\fill}%
\begin{longtable}[]{@{}
  >{\raggedright\arraybackslash}p{(\AFSlideTableWidth - 2\tabcolsep) * \real{0.6842}}
  >{\raggedright\arraybackslash}p{(\AFSlideTableWidth - 2\tabcolsep) * \real{0.3158}}@{}}
\label{tbl:nlevel3-params}\tabularnewline
\toprule\noalign{}
\begin{minipage}[b]{\linewidth}\raggedright
Parameter
\end{minipage} & \begin{minipage}[b]{\linewidth}\raggedright
Value
\end{minipage} \\
\midrule\noalign{}
\endfirsthead
\toprule\noalign{}
\begin{minipage}[b]{\linewidth}\raggedright
Parameter
\end{minipage} & \begin{minipage}[b]{\linewidth}\raggedright
Value
\end{minipage} \\
\midrule\noalign{}
\endhead
Agents & 4 \\
Trials & 960 \\
Acuity & 0.85 \\
\bottomrule\noalign{}
\end{longtable}
\endgroup
\end{frame}

\begin{frame}{Study parameters for the\ldots{} (part 2)}
The executed three-level configuration is summarized in
tbl.~18.
\end{frame}

\end{document}
